% =====================================================================
%  CHƯƠNG 2 — TỔNG QUAN
% =====================================================================
\chapter{Tổng quan}

\section{Bản chất rủi ro kiệt quệ tài chính}

Trong tài chính doanh nghiệp, kiệt quệ tài chính (financial distress) là trạng thái doanh nghiệp
không còn đủ khả năng tạo dòng tiền để thực hiện đầy đủ các nghĩa vụ nợ, hoặc đang xuất hiện các
dấu hiệu dẫn tới tình trạng đó \cite{beaver1966}. Các dấu hiệu thường được dùng gồm lỗ ròng kéo dài,
dòng tiền hoạt động âm, vốn lưu động âm và vốn chủ sở hữu âm.

Cần phân biệt rõ hai khái niệm. Kiệt quệ tài chính là một \emph{quá trình} kéo dài nhiều quý, còn
phá sản là một \emph{sự kiện} pháp lý xảy ra tại một thời điểm. Doanh nghiệp có thể ở trạng thái
kiệt quệ tài chính trong vài quý mà chưa nộp đơn phá sản, và ngược lại một số vụ phá sản xảy ra
đột ngột do nguyên nhân không phản ánh trong báo cáo tài chính quý. Khóa luận chọn định vị bài toán
ở mức kiệt quệ tài chính vì lý do thực nghiệm: kiểm tra hồ sơ 8-K cho thấy bộ dữ liệu hiện tại
không chứa sự kiện phá sản nào (0/8 công ty).

Với ngành bán lẻ, danh sách tín hiệu sớm cần chú ý khác với doanh nghiệp sản xuất. Vòng quay hàng
tồn và tỷ lệ chi phí bán hàng trên doanh thu phản ứng nhanh nhất khi sức mua suy yếu; cơ cấu nợ
ngắn hạn quyết định khả năng chịu được giai đoạn doanh thu đi ngang; và dòng tiền hoạt động là chỉ
báo ít bị pha loãng bởi các bút toán kế toán hơn lợi nhuận ròng.

\section{Chuẩn dữ liệu báo cáo tài chính SEC XBRL}

Từ năm 2009, U.S. Securities and Exchange Commission (SEC) yêu cầu công ty đại chúng nộp báo cáo
tài chính kèm phần XBRL (eXtensible Business Reporting Language) \cite{secxbrl2026}. Mỗi dòng số
liệu được gắn một thẻ (tag) thuộc taxonomy US-GAAP. SEC công bố các fact này qua endpoint
\texttt{data.sec.gov/api/xbrl/companyfacts/CIK\{cik\}.json} \cite{secfacts2026}. Mỗi fact có cấu
trúc gồm: \texttt{tag}, kỳ bắt đầu \texttt{start}, kỳ kết thúc \texttt{end}, giá trị \texttt{val},
số hiệu hồ sơ \texttt{accn} và ngày nộp \texttt{filed}.

\subsection{Hai loại chỉ tiêu cần phân biệt}

\begin{itemize}
\item Chỉ tiêu \textbf{số dư} (instant), ví dụ \texttt{Assets}, \texttt{Liabilities},
      \texttt{StockholdersEquity}: chỉ có \texttt{end}, đặc trưng cho một thời điểm.
\item Chỉ tiêu \textbf{số phát sinh} (duration), ví dụ \texttt{Revenues}, \texttt{NetIncomeLoss},
      \texttt{NetCashProvidedByOperatingActivities}: có cả \texttt{start} và \texttt{end}.
\end{itemize}

Với chỉ tiêu số phát sinh, một điểm dễ sai là lấy nhầm fact luỹ kế từ đầu năm (year-to-date) thay
vì fact của riêng một quý. Nếu một công ty chỉ công bố dạng luỹ kế, phải trừ luỹ kế hiện tại cho
luỹ kế của kỳ liền trước để ra giá trị quý riêng. Một quý tài chính 13 tuần tương ứng khoảng 70--105
ngày; đây là cận lọc hữu ích để loại các fact bán niên hoặc luỹ kế.

Đơn vị tiền trong bộ dữ liệu được quy đổi minh hoạ sang VND chỉ để thống nhất đơn vị trình bày. Cần
lưu ý đây là dữ liệu theo chuẩn US-GAAP của Mỹ, không phải báo cáo tài chính Việt Nam. Khi đối chiếu
với bối cảnh trong nước phải tính đến khác biệt về chế độ kế toán doanh nghiệp \cite{thongtu200},
các chuẩn mực kế toán Việt Nam \cite{vas2005} và quy định pháp lý về kế toán \cite{luatketoan2015}.

\section{Phương pháp dự báo cổ điển so với học máy}

\subsection{Altman Z-score}

Altman xây dựng điểm số tổ hợp tuyến tính từ vài tỷ số kế toán với trọng số cố định
\cite{altman1968}. Với doanh nghiệp phi sản xuất và dịch vụ, biến thể $Z''$ được dùng:

\begin{equation}
Z'' = 6{,}56\,\frac{\text{vốn lưu động}}{\text{tổng tài sản}}
    + 3{,}26\,\frac{\text{lợi nhuận giữ lại}}{\text{tổng tài sản}}
    + 6{,}72\,\frac{\text{EBIT}}{\text{tổng tài sản}}
    + 1{,}05\,\frac{\text{vốn chủ sở hữu}}{\text{tổng nợ}}
\end{equation}

Ưu điểm của $Z''$ là minh bạch và không cần huấn luyện; hạn chế là trọng số được chọn từ một tập
mẫu cũ, giả định quan hệ tuyến tính giữa các tỷ số và dùng một ngưỡng chung cho mọi ngành.
Barboza và cộng sự cho thấy các mô hình học máy vượt các mô hình thống kê truyền thống trên bài
toán dự báo phá sản \cite{barboza2017}.

\subsection{Học máy}

Học máy thay hai giả định của mô hình cổ điển — trọng số cố định và quan hệ tuyến tính — bằng cách
học trọng số và mặt quyết định từ dữ liệu. Với 47 đặc trưng và các tương tác phi tuyến (ví dụ lợi
nhuận giảm chỉ nguy hiểm khi đi kèm đòn bẩy cao), mô hình cây và mạng nơ-ron có cơ hội bắt được
quan hệ mà một công thức tuyến tính bỏ lỡ. Khóa luận chọn bốn họ đại diện cho bốn cơ chế học khác
nhau: tuyến tính (Logistic Regression) \cite{pedregosa2011}, bagging cây (Random Forest)
\cite{breiman2001}, boosting cây (HistGradientBoosting) \cite{friedman2001} và mạng nơ-ron
(MLP). XGBoost được dùng như đối chứng phản biện \cite{chen2016}.

\section{Các thước đo đánh giá trên dữ liệu mất cân bằng}

\subsection{Vì sao không dùng độ chính xác}

Trên tập test của khóa luận, lớp suy giảm chiếm 59,4\%; chỉ cần dự đoán tất cả là suy giảm đã đạt
độ chính xác 59,4\% mà không phân biệt được gì. Vì vậy độ chính xác (accuracy) không được dùng làm
chỉ số chính. Thay vào đó khóa luận báo cáo AUROC \cite{fawcett2006}, AP (diện tích dưới đường
Precision--Recall) \cite{davis2006}, Precision/Recall/F1, macro-F1, MCC và Brier score.

\subsection{Ngưỡng quyết định theo chi phí kỳ vọng}

Ngưỡng 0,5 là mặc định nhưng không có cơ sở. Khóa luận chọn ngưỡng tối thiểu hoá chi phí kỳ vọng
$C_{FN}\cdot FN + C_{FP}\cdot FP$ (lý thuyết quyết định). Nếu xác suất được hiệu chuẩn đúng, ngưỡng
tối ưu là
\begin{equation}
p^{*} = \frac{C_{FP}}{C_{FN} + C_{FP}}.
\end{equation}
Với $C_{FN} = 5$ và $C_{FP} = 1$, $p^{*} = 1/6 \approx 0{,}167$. Trên dữ liệu thực, ngưỡng tối ưu
tính được cao hơn nhiều (0,788) vì mô hình không hiệu chuẩn tuyệt đối — đây cũng là một điểm cần
thảo luận khi triển khai.
